% !TEX root = userguide.tex

\def\headdate{28th March 2023}

\section{Flow of a generalized Newtonian fluid model}
\label{sec:gennewtonian}

In this section some examples of the flow of a generalized Newtonian fluid model
will be given. 

\subsection{A generalized Newtonian fluid problem on a unit square with periodic boundary
         conditions: channel problem with specified flowrate; collocation;
         Picard and Newton-Raphson iteration}
\label{sec:gennewtonian1and2}
                                                                                
We consider a generalized Newtonian fluid for a channel problem on a unit square
similar to the problem \texttt{stokes2.f90} in Sec.~\ref{sec:stokes2}. The
viscosity is given by a viscosity model (power law, Carreau or Carreau-Yasuda).
                                                                                
The example files are \texttt{generalized\_newtonian1.f90} and \texttt{generalized\_newtonian2.f90}
and can be obtained by typing at the
command line:
\begin{quote}
   \texttt{getexample generalized\_newtonian}
\end{quote}
In \texttt{generalized\_newtonian1.f90} the problem is solved
using a Picard iteration scheme. If the weak form is written as:
Find $\vek u$ and $p$ such that
\begin{alignat*}{3}
  &\big(\ten D_v,2\eta(\dot\gamma(\vek u))\ten D(u)\big)
    -(\nabla\cdot\vek v,p)&&=(\vek v,\vek t_{\text{N}})_{\Gamma_{\text{N}}}+(\vek v,\vek f)
&\qquad&\text{for all }\vek v\\
  &\,(q,\nabla\cdot\vek u)&& =0,&\qquad&\text{for all }q
\end{alignat*}
where $\dot\gamma(\vek u)=\sqrt{2\ten D(\vek u):\ten D(\vek u)}$, then the Picard iteration scheme can be written as
\begin{alignat*}{3}
  &\big(\ten D_v,2\eta(\dot\gamma(\vek u_{j}))\ten D(u_{j+1})\big)
    -(\nabla\cdot\vek v,p_{j+1})&&=(\vek v,\vek t_{\text{N}})_{\Gamma_{\text{N}}}+(\vek v,\vek f)
&\qquad&\text{for all }\vek v\\
  &\,(q,\nabla\cdot\vek u_{j+1})&& =0,&\qquad&\text{for all }q
\end{alignat*}
with $j=1,2,\ldots$ and initial guess $\vek u_0$.
The resulting velocity profile for a Carreau model has been plotted in
Fig.~\ref{fig:vprofile2}.
\begin{figure}[htbp!]
   \begin{center}
   \includegraphics[scale=0.55]{vprofile2}
   \end{center}
   \caption{Velocity profile in the flow of generalized Newtonian fluid problem. Carreau model
($\eta_0=1$, $\eta_\infty=0$, $n=0.3$, $\lambda=1$). Flow rate is 2.}
    \label{fig:vprofile2}
\end{figure}

In \texttt{generalized\_newtonian2.f90} the problem is solved
using a Newton-Raphson iteration scheme with a possibility to start with a number Picard iterations to
obtain better convergence. If the weak form is written as:
Find $\vek u$ and $p$ such that
\begin{alignat*}{3}
  &\big(\ten D_v,2\eta(\dot\gamma(\vek u))\ten D(\vek u)\big)
    -(\nabla\cdot\vek v,p)&&=(\vek v,\vek t_{\text{N}})_{\Gamma_{\text{N}}}+(\vek v,\vek f)
&\qquad&\text{for all }\vek v\\
  &\,(q,\nabla\cdot\vek u)&& =0,&\qquad&\text{for all }q
\end{alignat*}
where $\dot\gamma(\vek u)=\sqrt{2\ten D(\vek u):\ten D(\vek u)}$, then the Newton-Raphson iteration scheme can be written as
\begin{multline}
  \big(\ten D_v,\frac{4\eta'_j}{\dot\gamma_j}\ten D_j\ten D_j:\ten D(\Delta \vek u_{j+1})+
  2\eta_{j}\ten D(\Delta \vek u_{j+1})\big)-(\nabla\cdot\vek v,p_{j+1}) \\
= -(\ten D_v, 2 \eta_j \ten D_j)+(\vek v,\vek t_{\text{N}})_{\Gamma_{\text{N}}}+(\vek v,\vek f)
\end{multline}
\begin{equation}
  (q,\nabla\cdot\Delta \vek u_{j+1}) =-(q,\nabla\cdot\vek u_{j}),
\end{equation}
and
\begin{equation}
  \vek u_{j+1} = \vek u_j+\Delta \vek u_{j+1}
\end{equation}
with $j=1,2,\ldots$ and initial guess $\vek u_0$. Here, $\eta_j=\eta(\dot\gamma(\vek u_j))$, $\eta'_j=\lderiv{\eta}{\dot\gamma}$ evaluated for $\vek u_j$ and
 $\ten D_j=\ten D(\vek u_j)$.
 
In \texttt{generalized\_newtonian3.f90} the same problem is solved as in \texttt{generalized\_stokes4.f90} (see Sec.~\ref{sec:genstokes3}), i.e.\
the viscosity is given by a regularized Bingham model \cite{Papanastasiou87} and solved using Picard iteration. However, now the generalized Newtonian elements are used instead. In \texttt{generalized\_newtonian4.f90} the problem is extended to using Newton-Raphson iteration preceded by a number of Picard iterations. 

\subsection{A generalized Newtonian fluid problem with yield of the flow around a sphere in a cylindrical tube.
         Picard and Newton-Raphson iteration.}
\label{sec:gennewtoniansphere}
The example files are \texttt{sphere34.f90} and \texttt{sphere35.f90}
and can be obtained by typing at the
command line:
\begin{quote}
   \texttt{getexample sphere}
\end{quote}
The viscosity is given by a regularized yield model \cite{Papanastasiou87} and the problems are solved using 
Newton-Raphson iteration preceded by a number of Picard iterations. In \texttt{sphere34.f90} the velocity of sphere is imposed whereas in \texttt{sphere35.f90} the external force on the sphere is imposed using a constraint.
         

